\documentclass[12pt]{article}
\usepackage{latexblog}

% ============================================================
% Blog Metadata
% ============================================================
\blogtitle{Redis实现分布式锁}
\blogdate{2020-07-17}
\blogtags{Redis, 闲话编程, 存储}
\bloglang{zh}
\blogsummary{}

\begin{document}
\makeblogtitle

Redis一个比较重要的应用场景就是分布式锁DLM (Distributed Lock Manager)。实际上已经有很多现成的redis库来完成这个功能了，但是可能实现途径有所差别，那么，正确的做法是什么呢？Redis官方建议了一个算法叫做\texttt{Redlock}，可以将其作为起点去实现更复杂的方案，来研究一下它的思路。

\section{原则}\label{ux539fux5219}

要实现分布式锁得满足一些必要的条件：

\begin{itemize}
\tightlist
\item
  互斥性：任意时刻只能有一个客户端能够成功获取锁
\item
  避免死锁：即使锁的持有者crash
\item
  容错性：只要redis集群大多数节点正常，客户端就能够获取或者释放锁
\end{itemize}

有一些常见的做法并不是安全的锁实现方式。

最简单的做法是，获取锁时，在redis创建一个带有过期时间的key。当需要释放锁时，删除这个key。这种做法无法避免单点故障，如果redis master宕机，则无法成功获取或者释放锁了。如果添加一个slave节点呢？很不幸也行不通，因为redis的主从复制是异步的，由此带来竞争：

\begin{itemize}
\tightlist
\item
  Client A 在master节点上获取了锁
\item
  master节点在将数据同步到slave之前crash掉了
\item
  slave提升为新的master
\item
  Client B于是可以获取到相同的锁了（不满足互斥性）
\end{itemize}

\section{redis单机下的正确做法}\label{redisux5355ux673aux4e0bux7684ux6b63ux786eux505aux6cd5}

在不考虑redis集群的情况下，如何正确的实现一个分布式锁呢？其实也比较简单：

加锁：

\begin{Shaded}
\begin{Highlighting}[]
\ExtensionTok{SET}\NormalTok{ resource\_name my\_random\_value NX PX 30000}
\end{Highlighting}
\end{Shaded}

其中，\texttt{NX}保证只有在key不存在的情况下才会被设置到redis中，\texttt{PX\ 30000} 设置了过期时间为30000毫秒。而key的值被设置为一个随机值，这个值必须在所有的客户端和加锁请求中唯一。之所以要这么做，是为了保证能够安全的释放锁，只有当key存在，且是由锁的持有者发起的解锁请求的时候，才删除这个key：

\begin{Shaded}
\begin{Highlighting}[]
\ControlFlowTok{if} \VariableTok{redis}\OperatorTok{.}\NormalTok{call}\OperatorTok{(}\StringTok{"get"}\OperatorTok{,}\ConstantTok{KEYS}\OperatorTok{[}\DecValTok{1}\OperatorTok{])} \OperatorTok{==} \ConstantTok{ARGV}\OperatorTok{[}\DecValTok{1}\OperatorTok{]} \ControlFlowTok{then}
    \ControlFlowTok{return} \VariableTok{redis}\OperatorTok{.}\NormalTok{call}\OperatorTok{(}\StringTok{"del"}\OperatorTok{,}\ConstantTok{KEYS}\OperatorTok{[}\DecValTok{1}\OperatorTok{])}
\ControlFlowTok{else}
    \ControlFlowTok{return} \DecValTok{0}
\ControlFlowTok{end}
\end{Highlighting}
\end{Shaded}

主要避免的一个问题就是，锁被其他的客户端给错误地释放了（有可能客户端释放锁的时候，因为某些原因锁已经过期了，但是其他的客户端已经获得了锁）。而锁的过期时间（或者说有效期），应该足够client完成操作，避免任务在进行的过程中其他客户端又获得了锁。

在单机的情况下以上就实现了一个较为完美的锁，如果要扩展到redis集群呢？

\section{Redlock 算法}\label{redlock-ux7b97ux6cd5}

(TBD)

\begin{itemize}
\tightlist
\item
  \href{https://redis.io/topics/distlock}{Distributed locks with Redis}
\end{itemize}


\end{document}
